\subsection{Some generality}\label{subsec:some_generality}

The cohomology $H^n(T^2; \Z)$ of the torus is well known, so that nothing remains to be proven in the case of \textsf{p1}.

For the point group $P$ of any $2$-dimensional space group, the vanishing $H^3(T^2; \Z) = 0$ implies $E^{0, 3}_\infty = 0$, so that
\begin{gather*}
H^3_P\big(T^2; \Z\big) = F^0H^3_P\big(T^2; \Z\big) = F^1H^3_P\big(T^2; \Z\big).
\end{gather*}
Note that each point group $P$ f\/ixes a point on $T^2$, so that
\begin{gather*}
F^3H^3_P\big(T^2; \Z\big) = H^3_P(\pt; \Z)= H^3_{\mathrm{group}}(P; \Z)= H^2_{\mathrm{group}}(P; U(1)).
\end{gather*}
Then the main task for the proof of Theorem \ref{thm:main} is to compute $H^3_P(T^2; \Z)$ and $F^2H^3_P(T^2; \Z)$, since in the case where $P$ is the cyclic group $\Z_n$ or the dihedral group $D_n$, the cohomology $H^m_P(\pt; \Z)$ is summarized as follows:
$$
\begin{array}{|c|c|c|c|c|}
\hline
P & H^0_P(\pt; \Z) & H^1_P(\pt; \Z) & H^2_P(\pt; \Z) & H^3_P(\pt; \Z) \tsep{2pt}\bsep{2pt}\\
\hline
\Z_n & \Z & 0 & \Z_n & 0 \\
\hline
D_n & \Z & 0 &
\begin{cases}
\Z_2 & (\mbox{$n$: odd}) \\
\Z_2^{\oplus 2} & (\mbox{$n$: even})
\end{cases} &
\begin{cases}
0 & (\mbox{$n$: odd}) \\
\Z_2 & (\mbox{$n$: even})
\end{cases}\tsep{10pt}\bsep{10pt}\\
\hline
\end{array}
$$
The degree $0$ part $H^0_P(\pt; \Z) = H^0(BP; \Z) = \Z$ is clear. Since $P$ is f\/inite, the degree $1$ part $H^1_P(\pt; \Z) \cong \operatorname{Hom}(P, \Z)$ is trivial. The degree $2$ part $H^2_P(\pt; \Z) \cong \operatorname{Hom}(P, U(1))$ can be seen by the classif\/ication of irreducible representations. Finally, the degree $3$ part $H^3_P(\pt; \Z) \cong H^2_{\mathrm{group}}(P; U(1))$ for $P = \Z_n, D_n$ can be found in \cite{Ka}.

In the rest of the section, we may use a structure of $T^2$ as a $P$-CW complex. In general, for a compact Lie group $G$, a \textit{$G$-CW complex} is an analogue of a CW complex made of \textit{$G$-cells}. A~$d$-dimensional $G$-cell is a $G$-space of the form $G/H \times e^d$, where $H \subset G$ is a closed subgroup and~$e^d$ is the standard $d$-dimensional cell. The $G$-action on $G/H$ is the left translation, whereas that on~$e^d$ is trivial. For the details, we refer the reader to \cite{May}.

We later compute a group cohomology via cohomology of a space:

\begin{lem} \label{lem:1_dim_complex}
Suppose that a finite group $G$ acts on a path connected space $Y$ fixing at least one point $\pt \in Y$. Suppose further that $Y$ is a CW complex consisting of only cells of dimension less than or equal to $1$. Then the following holds true for all $n \ge 0$.
\begin{gather*}
H^n_{\mathrm{group}}\big(G; H^1(Y; \Z)\big) \cong \tilde{H}^{n+1}_G(Y; \Z),
\end{gather*}
where $\tilde{H}^{n+1}_G(Y; \Z)$ stands for the reduced cohomology.
\end{lem}

Notice that a $G$-CW complex is naturally a CW complex.

\begin{proof}
Consider the Leray--Serre spectral sequence
\begin{gather*}
E_2^{p, q} = H^p_{\mathrm{group}}(G; H^q(Y; \Z))\Longrightarrow H^*_G(Y; \Z).
\end{gather*}
Note that $H^q(Y; \Z) = 0$ for $q \neq 0, 1$. The $E_2$-term $E_2^{p, 0} = H^p_{\mathrm{group}}(G; \Z)$ must survive into the direct summand $H^p_G(\pt; \Z)$ in $H^p_G(Y; \Z) \cong H^p_G(\pt; \Z) \oplus \tilde{H}^p_G(Y; \Z)$. Therefore it must hold that
\begin{gather*}
\tilde{H}^p_G(Y; \Z) \cong E^{p-1, 1}_\infty = E^{p-1, 1}_2\cong H^{p-1}_{\mathrm{group}}\big(G; H^1(Y; \Z)\big),
\end{gather*}
which completes the proof.
\end{proof}

We also prepare a simple lemma about group cohomology: Let $G$ be a f\/inite group, $c \colon G \to \Z_2 = \{ \pm 1 \}$ a surjective homomorphism, and $\tilde{\Z} = \Z_c$ the $G$-module such that its underlying group is $\Z$ and $G$ acts (from the right) by $m \mapsto m c(g)$. A typical example is a f\/inite subgroup $P \subset {\rm O}(2)$ such that $P \not\subset {\rm SO}(2)$ with $c$ the composition of the inclusion $P \to {\rm O}(2)$ and the determinant ${\rm O}(2) \to \Z_2$.

\begin{lem} \label{lem:group_coh_orientation_coeff}
Let $G$, $c$ and $\tilde{\Z}$ be as above. Then, $H^0_{\mathrm{group}}(G; \tilde{\Z}) = 0$ and $H^1_{\mathrm{group}}(G; \tilde{\Z}) \cong \Z_2$.
\end{lem}

\begin{proof}For any $n \in C^0_{\mathrm{group}}(G; \tilde{\Z}) = \Z$, its coboundary $\partial n\colon G \to \Z$ is $(\partial n)(g) = n(1 - c(g))$. Thus, the assumption that $c$ is surjective implies the vanishing of the $0$th cohomology. The inclusion $\operatorname{Ker}(c) \subset G$ induces an injection on $1$-cocycles
\begin{gather*}
Z^1_{\mathrm{group}}\big(G; \tilde{\Z}\big) \to Z^1_{\mathrm{group}}\big(\operatorname{Ker}(c); \tilde{\Z}\big) = \operatorname{Hom}(\operatorname{Ker}(c), \Z) = 0.
\end{gather*}
Thus, given a group $1$-cocycle $\phi \in Z^1_{\mathrm{group}}(G; \tilde{\Z})$, it holds that $\phi(g) = 0$ for all $g \in \operatorname{Ker}(c)$. If $g, h \not\in \operatorname{Ker}(c)$, then the cocycle condition $(\partial \phi)(g, h) = 0$ implies $\phi(g) = \phi(h)$. Therefore $\phi \colon G \to \Z$ is always of the form $\phi(g) = n (1 - c(g))/2$ for some $n \in \Z$. This provides the identif\/ication $Z^1_{\mathrm{group}}(G; \tilde{\Z}) \cong \Z$ as well as $B^1_{\mathrm{group}}(G; \tilde{\Z}) \cong 2\Z$, which completes the proof.
\end{proof}

In some cases, the computations of the Leray--Serre spectral sequence are similar, which we summarize as follows:

\begin{lem} \label{lem:typical_argument_non_orientation_preserving_case}
Let $G$ be a finite group acting on the torus $T^2$ such that:
\begin{itemize}\itemsep=0pt
\item there is a fixed point $\pt \in T^2$,

\item the $G$-action does not preserve the orientation of $T^2$.
\end{itemize}
Then the following holds true about the Leray--Serre spectral sequence:
\begin{itemize}\itemsep=0pt
\item[$(a)$] $F^2H^3_G(T^2; \Z) \cong E_2^{2, 1} \oplus E_2^{3, 0}$,

\item[$(b)$] $H^n_G(T^2; \Z) \cong \bigoplus_{p + q = n} E_2^{p, q}$ for $n \le 2$.
\end{itemize}
\end{lem}

\begin{proof}
In the Leray--Serre spectral sequence
$E_2^{p, q} = H^p_{\mathrm{group}}(G; H^q(T^2; \Z))$, the coef\/f\/icient in the group cohomology $H^0(T^2) \cong \Z$ is identif\/ied with the trivial $G$-module, and $H^2(T^2) \cong \Z$ with the $G$-module in Lemma \ref{lem:group_coh_orientation_coeff}. Then the relevant $E_2$-terms can be summarized as follows:
$$
\begin{array}{|c|c|c|c|c|c|}
\hline
q = 3 & 0 & 0 & 0 & 0 & 0 \\
\hline
q = 2 & 0 & \Z_2 & & & \\
\hline
q = 1 & E_2^{0, 1} & E_2^{1, 1} & E_2^{2, 1} & & \tsep{2pt}\bsep{2pt}\\
\hline
q = 0 & E_2^{0, 0} & E_2^{1, 0} & E_2^{2, 0} & E_2^{3, 0} & E_2^{4, 0} \tsep{2pt}\bsep{2pt}\\
\hline
E_2^{p, q} \tsep{2pt}\bsep{2pt} & p = 0 & p = 1 & p = 2 & p = 3 & p = 4\\
\hline
\end{array}
$$
Since $G$ f\/ixes $\pt \in T^2$, we have the decomposition $H^n_G(T^2; \Z) \cong H^n_G(\pt; \Z) \oplus \tilde{H}^n_G(T^2; \Z)$, where $\tilde{H}^n_G(T^2; \Z)$ is the reduced cohomology. Therefore the $E_2$-term $E_2^{n, 0} = H^n_{\mathrm{group}}(G; \Z) \cong H^n_G(\pt; \Z)$ must survive into the direct summand $H^n_G(\pt; \Z)$ in $H^n_G(T^2; \Z)$. This implies that $E_2^{n, 0} = E_\infty^{n, 0}$ is always a direct summand of the subgroups $F^pH^n_G(T^2; \Z) \subset H^n_G(T^2; \Z)$ and that $d_2 \colon E_2^{p-2, 1} \to E_2^{p, 0}$ is trivial. As a result, we get $E_2^{2, 1} = E_\infty^{2, 1}$ and the isomorphism (a). Also $E_2^{p, q} = E_\infty^{p, q}$ for $p + q \le 2$, and the isomorphism (b) follows.
\end{proof}
